\documentclass[a4paper,12pt]{article}
\usepackage[utf8]{inputenc}
\usepackage[T2A]{fontenc}
\usepackage[russian]{babel}
\usepackage{booktabs}
\usepackage{geometry}
\usepackage{siunitx}
\geometry{margin=2cm}
\sisetup{group-separator={\,}}

\title{Практикум №1 «Методы сортировки»\\Отчет по экспериментальному сравнению}
\author{Ларичева Дарья}
\date{}

\begin{document}
\maketitle

\section*{Постановка}
Реализованы две сортировки массива целых чисел по возрастанию: \textbf{Selection Sort} и \textbf{Shaker Sort}. Для каждого метода измерялись:
\begin{itemize}
  \item число сравнений элементов;
  \item число перемещений (присваиваний). Для обмена двух элементов учитывалось 3 перемещения.
\end{itemize}

Сравнение проводилось на одинаковых исходных массивах для каждого метода. Рассматривались типы входных данных:
\begin{enumerate}
  \item уже упорядочен;
  \item упорядочен в обратном порядке;
  \item случайный;
  \item случайный (другой массив).
\end{enumerate}
Длины массивов: $n \in \{10, 100, 1000, 10000\}$. Для каждого случая выполнялось 10 запусков, в таблицах приведены средние значения.

\section*{Результаты}
\begin{table}[h!]
\centering
\caption{Средние значения для $n=10$ (10 запусков)}
\begin{tabular}{l S[round-mode=places,round-precision=0] S[round-mode=places,round-precision=0] S[round-mode=places,round-precision=0] S[round-mode=places,round-precision=0]}
\toprule
{Тип массива} & {Sel: сравнения} & {Sel: перемещения} & {Shaker: сравнения} & {Shaker: перемещения} \\
\midrule
Упорядочен & 45 & 0 & 9 & 0 \\
Обратный & 45 & 15 & 45 & 135 \\
Случайный & 45 & 23 & 35 & 67 \\
Случайный 2 & 45 & 21 & 36 & 73 \\
\bottomrule
\end{tabular}
\end{table}

\begin{table}[h!]
\centering
\caption{Средние значения для $n=100$ (10 запусков)}
\begin{tabular}{l S[round-mode=places,round-precision=0] S[round-mode=places,round-precision=0] S[round-mode=places,round-precision=0] S[round-mode=places,round-precision=0]}
\toprule
{Тип массива} & {Sel: сравнения} & {Sel: перемещения} & {Shaker: сравнения} & {Shaker: перемещения} \\
\midrule
Упорядочен & 4950 & 0 & 99 & 0 \\
Обратный & 4950 & 150 & 4950 & 14850 \\
Случайный & 4950 & 288 & 3442 & 7522 \\
Случайный 2 & 4950 & 287 & 3314 & 7258 \\
\bottomrule
\end{tabular}
\end{table}

\begin{table}[h!]
\centering
\caption{Средние значения для $n=1000$ (10 запусков)}
\begin{tabular}{l S[round-mode=places,round-precision=0] S[round-mode=places,round-precision=0] S[round-mode=places,round-precision=0] S[round-mode=places,round-precision=0]}
\toprule
{Тип массива} & {Sel: сравнения} & {Sel: перемещения} & {Shaker: сравнения} & {Shaker: перемещения} \\
\midrule
Упорядочен & 499500 & 0 & 999 & 0 \\
Обратный & 499500 & 1500 & 499500 & 1498500 \\
Случайный & 499500 & 2979 & 336132 & 754527 \\
Случайный 2 & 499500 & 2976 & 331432 & 741169 \\
\bottomrule
\end{tabular}
\end{table}

\begin{table}[h!]
\centering
\caption{Средние значения для $n=10000$ (10 запусков)}
\begin{tabular}{l S[round-mode=places,round-precision=0] S[round-mode=places,round-precision=0] S[round-mode=places,round-precision=0] S[round-mode=places,round-precision=0]}
\toprule
{Тип массива} & {Sel: сравнения} & {Sel: перемещения} & {Shaker: сравнения} & {Shaker: перемещения} \\
\midrule
Упорядочен & 49995000 & 0 & 9999 & 0 \\
Обратный & 49995000 & 15000 & 49995000 & 149985000 \\
Случайный & 49995000 & 29968 & 33441166 & 75236716 \\
Случайный 2 & 49995000 & 29969 & 33364674 & 75111827 \\
\bottomrule
\end{tabular}
\end{table}

\section*{Вывод}
Selection Sort делает фиксированное количество сравнений порядка $O(n^2)$ почти независимо от исходного порядка. Shaker Sort на упорядоченном массиве работает существенно лучше (мало обменов), но в худшем случае также имеет квадратичную сложность.

\end{document}
